\section{Class PSPACE --- TQBFARM --- Semantics}
%%% generated by the blueprint pipeline from TCSlib.Complexity.ClassPSPACE.TQBFARM.Semantics
\subsection*{Class PSPACE --- TQBFARM --- Semantics: Introduction}

Declarations of the Lean module \texttt{TCSlib.Complexity.ClassPSPACE.TQBFARM.Semantics}.

\subsection*{Class PSPACE --- TQBFARM --- Semantics: Results}

\begin{definition}[Syntax automaton states]
\lean{Complexity.LogProg.TQBFARM.SynSt}
\label{Complexity.LogProg.TQBFARM.SynSt}
\leanok
The syntax automaton state space is the disjoint union of an optional Boolean value and
a six-element state set: either no Boolean value has yet been recorded, a Boolean value
has been recorded, or the automaton is in one of six finite states.
\end{definition}

\begin{definition}[Syntax automaton transition]
\lean{Complexity.LogProg.TQBFARM.synStep}
\label{Complexity.LogProg.TQBFARM.synStep}
\leanok
\uses{Complexity.LogProg.TQBFARM.SynSt, Std.Sat.CNF.syntaxStep}
Define a transition on syntax states for each input bit. From the optional-value
component, an unrecorded bit is stored; after a stored false, a false clears the stored
value and a true enters state $0$; after a stored true, a true clears it and a false
enters state $5$. From one of the six finite states, apply the six-state transition:
from state $0$, move to $1$ on $1$ and to $4$ on $0$; from state $1$, move to $2$ on $1$
and to $0$ on $0$; from state $2$, move to $2$ on $1$ and to $3$ on $0$; from state $3$,
move to $1$ on either bit; and from states $4$ or $5$, move to state $5$ on either bit.
\end{definition}

\begin{definition}[Syntax acceptance test]
\lean{Complexity.LogProg.TQBFARM.synOK}
\label{Complexity.LogProg.TQBFARM.synOK}
\leanok
\uses{Complexity.LogProg.TQBFARM.synStep}
For a finite bit string $x$, start the syntax automaton in the state with no stored
Boolean value and process the bits of $x$ in order. The result is true exactly when the
final state is state $4$.
\end{definition}

\begin{theorem}[Syntax scan state identity]
\lean{Complexity.LogProg.TQBFARM.synStep_matrix}
\label{Complexity.LogProg.TQBFARM.synStep_matrix}
\leanok
\uses{Complexity.LogProg.TQBFARM.synStep, Std.Sat.CNF.syntaxStep}
\difficulty{4}
For every finite Boolean list $x$ and every state $q \in \{0,1,2,3,4,5\}$, scanning $x$
from the combined syntax state corresponding to $q$ ends at the combined syntax state
corresponding to the result of scanning $x$ from $q$ using the six-state syntax
transition.
\end{theorem}

\begin{theorem}[Syntax Acceptance and CNF Parsing]
\lean{Complexity.LogProg.TQBFARM.synOK_spec}
\label{Complexity.LogProg.TQBFARM.synOK_spec}
\leanok
\uses{Complexity.LogProg.TQBFARM.synOK, Complexity.LogProg.TQBFARM.synStep, Complexity.LogProg.TQBFARM.synStep_matrix, Std.Sat.CNF.parse, Std.Sat.CNF.syntaxOK_spec, Std.Sat.CNF.syntaxStep, Std.Sat.CNF.syntaxSuffix_run, Turing.pairDecode}
\difficulty{5}
For every bit string $x$, the syntax automaton accepts $x$ exactly when the aligned pair
parser succeeds on $x$ and the second component returned by that parser is successfully
parsed as a CNF formula, consuming the entire component.
\end{theorem}

\begin{theorem}[Truth for an empty matrix]
\lean{Complexity.LogProg.TQBFARM.truthAux_nil}
\label{Complexity.LogProg.TQBFARM.truthAux_nil}
\leanok
\uses{Complexity.QBF.Quant, Complexity.QBF.truthAux}
\difficulty{4}
For every finite quantifier prefix $q$, Boolean assignment
$\sigma:\mathbb{N}\to\{\mathrm{false},\mathrm{true}\}$, and index $i\in\mathbb{N}$, the
quantified truth predicate for the empty CNF matrix holds.
\end{theorem}

\begin{theorem}[Truth when the syntax check fails]
\lean{Complexity.LogProg.TQBFARM.truth_of_not_synOK}
\label{Complexity.LogProg.TQBFARM.truth_of_not_synOK}
\leanok
\uses{Complexity.LogProg.TQBFARM.synOK, Complexity.LogProg.TQBFARM.synOK_spec, Complexity.LogProg.TQBFARM.truthAux_nil, Complexity.QBF.decode, Complexity.QBF.quantOfBit, Complexity.QBF.truth, Std.Sat.CNF.decode, Std.Sat.CNF.fallback, Std.Sat.CNF.parse, Turing.pairDecode}
\difficulty{4}
For any bit string $x$, if its syntax check returns false, then the quantified Boolean
formula decoded from $x$ is true.
\end{theorem}

\begin{theorem}[Accepted Syntax Has a CNF Representation]
\lean{Complexity.LogProg.TQBFARM.repr_of_synOK}
\label{Complexity.LogProg.TQBFARM.repr_of_synOK}
\leanok
\uses{Complexity.LogProg.TQBFARM.synOK, Complexity.LogProg.TQBFARM.synOK_spec, Std.Sat.CNF.parse, Std.Sat.CNF.parse_repr, Std.Sat.CNF.serialize, Turing.eq_pairEncode_of_pairDecode, Turing.pairDecode, Turing.pairEncode}
\difficulty{4}
For every bit string $x$, if the syntax automaton accepts $x$, then there are bit
strings $q$ and a CNF formula $\varphi$ over natural-number-indexed variables such that
$x=\langle q,\mathrm{enc}(\varphi)\rangle$, where the pairing is the self-delimiting
encoding and $\mathrm{enc}(\varphi)$ is the specified binary serialization of the
formula.
\end{theorem}

\begin{theorem}[Decoding an encoded quantified formula]
\lean{Complexity.LogProg.TQBFARM.decode_pairEncode}
\label{Complexity.LogProg.TQBFARM.decode_pairEncode}
\leanok
\uses{Complexity.QBF.decode, Complexity.QBF.quantOfBit, Std.Sat.CNF.decode_serialize, Std.Sat.CNF.serialize, Turing.pairDecode_pairEncode, Turing.pairEncode}
\difficulty{1}
For every finite Boolean word $q$ and every CNF formula $\varphi$ over variables indexed
by $\mathbb{N}$, decoding the pair encoding of $q$ and the binary serialization of
$\varphi$ gives the quantified Boolean formula with quantifier prefix obtained by
interpreting each true bit as existential and each false bit as universal, and matrix
$\varphi$.
\end{theorem}

\begin{definition}[Binary assignment from a stack]
\lean{Complexity.LogProg.TQBFARM.asg}
\label{Complexity.LogProg.TQBFARM.asg}
\leanok
For natural numbers $k$ and $S$, define an assignment to each natural-number variable
$v$ as follows: it is true exactly when $v<k$ and the bit of $S$ at position $k-1-v$ is
$1$; it is false otherwise.
\end{definition}

\begin{theorem}[The empty-depth assignment]
\lean{Complexity.LogProg.TQBFARM.asg_zero}
\label{Complexity.LogProg.TQBFARM.asg_zero}
\leanok
\uses{Complexity.LogProg.TQBFARM.asg}
\difficulty{1}
For every natural number $S$, the assignment function at depth $0$ is constantly false.
\end{theorem}

\begin{theorem}[Pushing a bit onto an assignment]
\lean{Complexity.LogProg.TQBFARM.asg_push}
\label{Complexity.LogProg.TQBFARM.asg_push}
\leanok
\uses{Complexity.LogProg.TQBFARM.asg}
\difficulty{6}
For natural numbers $k,S$ and a bit $b$, updating the assignment function at index $k$
to $b$ gives the assignment function for depth $k+1$ and stack value $2S+b$, where $b$
is identified with $0$ or $1$.
\end{theorem}

\begin{definition}[Truth of a quantified Boolean formula]
\lean{Complexity.LogProg.TQBFARM.Val}
\label{Complexity.LogProg.TQBFARM.Val}
\leanok
\uses{Complexity.LogProg.TQBFARM.asg, Complexity.QBF.quantOfBit, Complexity.QBF.truthAux}
Given a Boolean list $qs$, a conjunctive normal form formula $\varphi$ over
natural-number variables, and natural numbers $k,S$, the predicate holds exactly when
$\varphi$ is true under the quantifier prefix obtained by dropping the first $k$ bits of
$qs$ and interpreting the remaining bits as quantifiers, with variables assigned by the
binary assignment determined by $k$ and $S$ and with parameter $k$.
\end{definition}

\begin{theorem}[Evaluation at the quantifier-prefix leaf]
\lean{Complexity.LogProg.TQBFARM.val_leaf}
\label{Complexity.LogProg.TQBFARM.val_leaf}
\leanok
\uses{Complexity.LogProg.TQBFARM.Val, Complexity.LogProg.TQBFARM.asg, Complexity.QBF.truthAux}
\difficulty{1}
For every finite Boolean quantifier prefix $q$, CNF formula $\varphi$, and natural
number $S$, the valuation predicate at depth $|q|$ is equivalent to the matrix $\varphi$
evaluating to true under the assignment encoded by $S$ at that depth.
\end{theorem}

\begin{theorem}[Mapping a List After Dropping Its First Entry]
\lean{Complexity.LogProg.TQBFARM.drop_map_eq}
\label{Complexity.LogProg.TQBFARM.drop_map_eq}
\leanok
\uses{Complexity.QBF.quantOfBit}
\difficulty{1}
For a Boolean list $q$ and a natural number $k<|q|$, mapping the bit-to-quantifier
interpretation over the list obtained by dropping its first $k$ entries gives the
quantifier corresponding to $q_k$, followed by the mapped list obtained by dropping the
first $k+1$ entries.
\end{theorem}

\begin{theorem}[Existential quantifier branching]
\lean{Complexity.LogProg.TQBFARM.val_ex}
\label{Complexity.LogProg.TQBFARM.val_ex}
\leanok
\uses{Complexity.LogProg.TQBFARM.Val, Complexity.LogProg.TQBFARM.asg_push, Complexity.LogProg.TQBFARM.drop_map_eq, Complexity.QBF.quantOfBit, Complexity.QBF.truthAux}
\difficulty{4}
Let $q$ be a finite Boolean quantifier prefix, $\varphi$ a CNF formula, and
$k,S\in\mathbb{N}$. If $k$ is a valid index in $q$ and its $k$th bit is true, then the
valuation predicate at $(k,S)$ holds exactly when it holds at $(k+1,2S)$ or at
$(k+1,2S+1)$.
\end{theorem}

\begin{theorem}[Universal quantifier branching]
\lean{Complexity.LogProg.TQBFARM.val_all}
\label{Complexity.LogProg.TQBFARM.val_all}
\leanok
\uses{Complexity.LogProg.TQBFARM.Val, Complexity.LogProg.TQBFARM.asg_push, Complexity.LogProg.TQBFARM.drop_map_eq, Complexity.QBF.quantOfBit, Complexity.QBF.truthAux}
\difficulty{4}
Let $q$ be a finite Boolean quantifier prefix, $\varphi$ a CNF formula, and
$k,S\in\mathbb{N}$. If $k$ is a valid index in $q$ and its $k$th bit is false, then the
valuation predicate at $(k,S)$ holds exactly when it holds at both $(k+1,2S)$ and
$(k+1,2S+1)$.
\end{theorem}

\begin{theorem}[Valuation at the root encoding]
\lean{Complexity.LogProg.TQBFARM.val_root}
\label{Complexity.LogProg.TQBFARM.val_root}
\leanok
\uses{Complexity.LogProg.TQBFARM.Val, Complexity.LogProg.TQBFARM.asg_zero, Complexity.LogProg.TQBFARM.decode_pairEncode, Complexity.QBF.decode, Complexity.QBF.truth, Std.Sat.CNF.serialize, Turing.pairEncode}
\difficulty{3}
For every finite Boolean quantifier prefix $q$ and CNF formula $\varphi$, the valuation
predicate at depth zero with number $1$ holds exactly when the quantified Boolean
formula decoded from the pairing of $q$ with the serialization of $\varphi$ is true.
\end{theorem}
